\section{Marginal Hub Investment}

\subsection{Hub Increment Cost Components}

The transit layer cost is the \textit{marginal} cost of adding passenger facilities to a diamond hub that is already being constructed for freight purposes. This is the critical distinction from standalone bus terminal analysis: the access roads, site utilities, stormwater management, security infrastructure, and parking land are already funded by the diamond interchange program (Component 2, \$4.5B). The transit increment funds only the passenger-specific elements.

Per T1/T1 hub, the transit increment components and cost ranges are:

\begin{table}[h]
\centering
\caption{Transit Hub Increment Cost Components per T1/T1 Hub}
\label{tab:hubcost}
\begin{tabular}{lrr}
\toprule
Component & Low Estimate & High Estimate \\
\midrule
Passenger platform with weather protection (4--8 coach bays) & \$0.8M & \$1.2M \\
Passenger parking (200--400 spaces, paved, ADA) & \$0.4M & \$0.8M \\
Ticketing/information kiosks (2--4 units) & \$0.15M & \$0.25M \\
ADA-accessible restrooms (increment over rest area standard) & \$0.05M & \$0.15M \\
Wi-Fi and device charging stations & \$0.08M & \$0.12M \\
Safety lighting upgrade for passenger areas & \$0.15M & \$0.25M \\
Accessible pathways (parking to platform) & \$0.12M & \$0.23M \\
\midrule
\textbf{Total per T1/T1 hub} & \textbf{\$1.75M} & \textbf{\$3.0M} \\
\textbf{Average} & \multicolumn{2}{c}{\$2.1M} \\
\bottomrule
\end{tabular}
\end{table}

For T1/T2 regional stops, the increment is smaller: the stop serves 2--4 coach bays rather than 4--8, and passenger parking is sized at 50--150 spaces. Average T1/T2 stop increment: \$0.8M (range \$0.5M--\$1.2M).

Total hub program cost:
\begin{align*}
\text{9 T1/T1 hubs} &\times \$2.1\text{M average} = \$18.9\text{M} \\
\text{50 T1/T2 stops} &\times \$0.8\text{M average} = \$40.0\text{M} \\
\text{Subtotal platform/facility increment} &= \$58.9\text{M}
\end{align*}

The remaining \$1.94 billion of the \$2 billion transit layer budget covers: (1) access road improvements and signage for passenger routing (\$0.8B), (2) grid and utility upgrades required for DCFC charging and facility electrical loads (\$0.6B), (3) program management, environmental review, and ADA compliance certification (\$0.4B), and (4) a contingency reserve of 7\% (\$0.14B). The access road and utility components are the largest cost drivers; they are non-negotiable because the existing rest area access infrastructure was designed for trucks, not for passenger vehicle flow patterns at hub boarding times.

\subsection{Cost-Efficiency Comparison}

The fundamental comparison is between the transit hub increment (\$161 per annual traveler served) and standalone bus terminal construction (\$2,200 per annual traveler served). Three reference points anchor the standalone estimate:

\begin{enumerate}
  \item \textbf{Raleigh, NC Union Station (2024)}: \$95 million capital for an estimated 420,000 annual intercity bus and rail passengers at opening. Capital cost per annual traveler: \$226. This is an urban hub in a high-demand corridor; the per-traveler cost is lower than the national average because the Raleigh corridor (NC--Northeast) has strong baseline demand \citep{Amtrak2023}.
  \item \textbf{Detroit Greyhound terminal replacement (2018)}: \$32 million for approximately 180,000 annual passengers. Capital cost per annual traveler: \$178. Detroit is a legacy Greyhound hub with established demand; the comparison understates the per-traveler cost for new markets \citep{BTS_intermodal2023}.
  \item \textbf{Average small-city bus terminal (estimated)}: For a standalone bus terminal in a secondary market (population 50,000--200,000) serving an estimated 30,000--50,000 annual passengers at full operation, the construction cost range is \$20--50 million, producing a cost-per-traveler of \$500--1,667. The \$2,200 estimate used as the comparison baseline reflects the higher end of this range, appropriate for markets where the hub must be built before demand is proven.
\end{enumerate}

The T1 hub estimate of \$161 per annual traveler is calculated as:
\begin{equation}
\text{Cost per traveler} = \frac{\$2.1\text{M average hub cost}}{13{,}100 \text{ estimated annual passengers per hub}} = \$160.30 \approx \$161
\end{equation}

The 13,100 annual passenger estimate per hub is derived from the hub transit-shed population (average 2,900,000 within 30 miles across all 9 hubs, non-overlapping), with a 0.45\% annual trip rate for transit-dependent households, calibrated to current Greyhound ridership density on comparable corridor segments.

\subsection{Comparison with Amtrak Capital Investment}

Amtrak's federal capital subsidy per annual rider has averaged approximately \$2,800 in recent years, based on FY2023 capital appropriations of \$2.1 billion and approximately 750,000 annual long-distance riders served by capital-funded equipment and infrastructure \citep{Amtrak2023}. The T1 hub investment at \$161 per traveler served is approximately 17 times more efficient than Amtrak capital investment per rider.

This comparison is not an argument against Amtrak; it is an argument for the transit hub model as a complementary intercity mobility investment, particularly in the corridors and markets where Amtrak does not and will not operate. The Southeast, Gulf Coast, and Great Plains have minimal Amtrak service and the highest concentrations of transit-dependent populations; the T1 hub network is the only realistic intercity transit investment in these markets at current federal budget levels \citep{IIJA2021}.
